% ৫. ডেটাবেজ তৈরি
\section{ডেটাবেজ তৈরি}
ভালো ডেটাবেজ তৈরির আগে কী তথ্য রাখা হবে, কোন কোন টেবিল লাগবে, কোন ফিল্ডের ডেটা টাইপ কী হবে এবং কোন ফিল্ডটি রেকর্ডকে অনন্যভাবে শনাক্ত করবে তা ঠিক করতে হয়।

\subsection{তৈরির ধাপ}
\begin{keypointbox}{ডেটাবেজ তৈরির ধাপ}
\begin{enumerate}
    \item উদ্দেশ্য নির্ধারণ: ডেটাবেজটি কোন কাজ করবে।
    \item এন্টিটি নির্বাচন: Student, Teacher, Product, Sales ইত্যাদি।
    \item ফিল্ড নির্ধারণ: Roll, Name, Class, Mobile, GPA ইত্যাদি।
    \item ডেটা টাইপ নির্ধারণ: Text, Number, Date/Time ইত্যাদি।
    \item Primary key নির্বাচন: যেমন StudentID বা Roll।
    \item টেবিল তৈরি, ডেটা এন্ট্রি, সম্পর্ক তৈরি এবং পরীক্ষা।
\end{enumerate}
\end{keypointbox}

\subsection{Table ও Field ডিজাইন}
\begin{examplebox}{Student ডেটাবেজ ডিজাইন}
\begin{center}
\begin{tabular}{|l|l|l|}
\hline
Field Name & Data Type & Note \\
\hline
StudentID & Number/AutoNumber & Primary Key \\
Name & Short Text & ছাত্রের নাম \\
GroupName & Short Text & বিভাগ \\
BirthDate & Date/Time & জন্মতারিখ \\
Mobile & Short Text & শুরুতে 0 থাকায় Text ভালো \\
\hline
\end{tabular}
\end{center}
\end{examplebox}

\begin{mistakebox}{সাধারণ ভুল}
মোবাইল নম্বর বা রোল নম্বর সব সময় গণনার জন্য ব্যবহৃত হয় না। তাই অনেক ক্ষেত্রে এগুলো Number না রেখে Text রাখা নিরাপদ, বিশেষ করে শুরুতে 0 থাকতে পারে এমন ডেটার ক্ষেত্রে।
\end{mistakebox}

\begin{practicebox}
একটি লাইব্রেরি ডেটাবেজের জন্য Book, Member ও Issue টেবিলের সম্ভাব্য ফিল্ড এবং Primary key নির্ধারণ করো।
\end{practicebox}
